\begin{frame}{VAE에는 이 항이 있다: 핵심 구조적 차이}
  \begin{columns}[T]
    \begin{column}{0.5\textwidth}
      \kb{VAE (15.1절)}
      \begin{itemize}
        \item KL 항 $D_{KL}(q_\phi(z|x) \| p(z))$이 ``인코더 출력
          분포가 표준정규분포에 가깝게 유지''를 \textbf{목적함수 안에
          명시적으로} 요구
        \item 잠재 공간 \textbf{전체}가 쓰이는 구조적 압력
        \item ``잠재 공간의 일부만 쓰는'' 퇴화 = \textbf{ELBO를
          직접적으로 낮추는 행위}
      \end{itemize}
    \end{column}
    \begin{column}{0.5\textwidth}
      \kb{GAN}
      \begin{itemize}
        \item 내쉬 균형 $p_G = p_{\text{data}}$ (모든 모드를 덮는
          분포)가 ``완전한 균형''
        \item 실제 학습이 풀고 있는 \textbf{근사적} 게임에서는
          ``일부 모드를 덮고 $D$를 속이는'' 분포도 \textbf{국소적으로
          균형처럼} 작용
        \item 이 대응 \textbf{항이 없음}
      \end{itemize}
    \end{column}
  \end{columns}
  \vskip0.8em
  \kq{이것이 VAE/GAN 비교에서 반복되는 \textbf{핵심 구조적 차이}.}
\end{frame}
